\subsection{Sprint 2}

Na Sprint 2, a prioridade foi fazer o projeto avançar de forma mais organizada, principalmente em três frentes: validar a finalização da modelagem do banco de dados, melhorar a comunicação do time e me manter disponível para desimpedir colegas de AGES menores sempre que fosse necessário. Depois da experiência da sprint anterior, ficou claro que não bastava apenas criar tarefas ou cobrar entregas e que era preciso criar uma forma mais saudável e objetiva de acompanhar o trabalho.

Uma das minhas principais atividades foi verificar a nova modelagem do banco de dados. Analisei, junto de outros colegas AGES IV, se a estrutura proposta fazia sentido com os documentos enviados pelo cliente e se ela estava coerente com o que o produto precisava representar. Depois dessa revisão, aprovei a modelagem, pois ela passou a refletir melhor as informações esperadas para o projeto e dava mais segurança para que o desenvolvimento seguisse sem retrabalho grande nessa parte. Essa validação foi importante porque o banco servia como base para várias decisões do sistema, especialmente no cadastro e organização das edificações e conteúdos relacionados, podendo assim desimpedir o desenvolvimento do \textit{backend}.

Também participei da reorganização do time em \textit{squads}. Essa mudança surgiu como uma tentativa de reduzir a sensação de que as informações ficavam espalhadas ou presas em grupos menores. Com as \textit{squads}, cada pessoa passou a ter referências mais claras de quem procurar, o que facilitou a troca assíncrona e ajudou o time a entregar mais do que vinha entregando anteriormente. Para mim, essa foi uma mudança positiva, porque deixou a gestão menos centralizada e tornou mais fácil identificar bloqueios antes que eles virassem problemas maiores.

Além da parte de acompanhamento, fiz duas \textit{spikes} técnicas: uma sobre o uso de ícones com fotos e outra sobre o limitador do bairro do Centro Histórico. Essas investigações ajudaram a entender melhor possibilidades e limitações antes de transformar as ideias em implementação direta. Mesmo não sendo entregas finais por si só, as \textit{spikes} foram importantes para reduzir incertezas técnicas e apoiar decisões futuras do produto.

O problema dessa sprint foi a comunicação pouco amigável entre alguns integrantes de AGES III com os AGES IV. Esse conflito não foi resolvido da melhor forma durante a execução, e a situação acabou aparecendo com mais força na retrospectiva. Apesar de não ter sido uma solução ideal, a discussão serviu para deixar claro que a forma como o time se comunica impacta diretamente o andamento das tarefas e o ambiente de trabalho. Foi uma experiência desconfortável, mas também necessária para perceber que maturidade técnica não resolve tudo se a comunicação não acompanha.

A maior lição que levo da Sprint 2 é que \textit{squads} ajudam bastante quando o time precisa lidar com comunicação assíncrona. Ter pessoas de referência facilita pedir ajuda, entender responsabilidades e saber quem está à frente de cada parte. Além disso, ficou como aprendizado perceber quando uma retrospectiva precisa ser interrompida, principalmente quando a discussão deixa de tratar o problema e passa a se concentrar nas pessoas. Para os próximos passos, eu deveria continuar monitorando AGES III, II e I para garantir as entregas no prazo, além de verificar se a forma de organizar e gerenciar as \textit{squads} ainda fazia sentido para todos.
